\section{Mechanistic Interpretability: la promesa y los límites de la evidencia interna}

La behavioral evaluation solo puede observar entradas y salidas; la \term{mechanistic interpretability} (interpretabilidad mecanicista) intenta comprender el cómputo interno del modelo a partir de activations, features y circuits. El trabajo de dictionary-learning de Anthropic descompone la representación mixta de cada neuron en features más dispersas, lo que ofrece una herramienta de investigación para saber «cuándo se activa cierto concepto y cómo influye en las etapas posteriores».

\subsection{Feature, circuit hypothesis e intervention}

Una \term{feature} es un patrón interpretable extraído de activations de alta dimensión y no necesariamente corresponde a una sola neuron; un \term{circuit} es una ruta de cómputo formada por la interaction de varias features, attention y MLP. El proceso de investigación va de la activation collection y la feature discovery a la behavior correlation, y después verifica la causalidad mediante ablation, steering o counterfactual intervention.

\lecturefigure{14-interpretability-loop.png}{La interpretabilidad mecanicista extrae features a partir de las activations, formula una circuit hypothesis y la valida mediante intervention.}{Anthropic monosemanticity research; subtítulos locales 38:17--39:03; concepto redibujado.}

\begin{knowledgebox}{Lectura de la figura: una feature legible no equivale a «leer la mente» por completo}
La feature label proviene de ejemplos de activación y de la interpretación humana, por lo que puede pasar por alto polysemy, distributed representation o context dependency; una circuit hypothesis requiere causal intervention, y una explicación local tampoco demuestra la seguridad global. La interpretabilidad funciona mejor como evidencia complementaria de la evaluation, del red team y de la anomaly investigation que como release gate por sí sola.
\end{knowledgebox}

\begin{warningbox}{No confundir una anthropomorphic label con la intención del modelo}
Etiquetas como «engaño», «poder» o «autoprotección» ayudan a localizar patrones de conducta, pero la feature activation interna del modelo no implica automáticamente un objetivo persistente ni una intención subjetiva. El informe debe indicar la operational definition, el dataset, los false positives, el intervention result y las alternative explanations.
\end{warningbox}

\subsection{Resumen de la sección}

La mechanistic interpretability aporta evidence interna, pero todavía se encuentra en una etapa temprana de investigación. La feature discovery, la causal validation y la behavioral evaluation deben validarse entre sí; no se puede tomar una gráfica de activations como prueba de seguridad.
